% !TEX root = ../mainfile.tex
% !TEX encoding = UTF-8 Unicode
% !TEX spellcheck = en_US

\chapter{Time-Varying Correlations}
\label{ch:dcc_garch}

% Target: 5-6 pages


\section{Univariate GARCH Estimates}
\label{sec:dcc_univariate}

% TODO: Table — GARCH(1,1) parameter estimates for each series
% (omega, alpha, beta, persistence = alpha + beta)
% TODO: Brief discussion of volatility persistence and clustering


\section{Dynamic Conditional Correlations}
\label{sec:dcc_estimates}

% TODO: DCC parameter estimates (alpha_DCC, beta_DCC)
% TODO: Figure — rho(USD, VIX) over time with event dates shaded
% Key test: does this turn negative around Liberation Day?
% TODO: Figure — rho(Gold, VIX) over time
% Should remain positive throughout → gold as consistent safe haven
% TODO: Figure — rho(Oil, EM currencies) over time
% Should spike during Hormuz crisis
% TODO: Figure — rho(Oil exporters, Oil) over time
% Should strengthen during Iran shock


\section{Safe-Haven Regression Results}
\label{sec:dcc_safehaven_results}

% TODO: Table — Baur & Lucey (2010) regression results
% rho(i,j,t) = c + beta_1 D_tariff + beta_2 D_Iran_Jun + beta_3 D_Iran_Feb + eps
% TODO: Key results:
% - USD-VIX: negative beta during tariff event → safe-haven breakdown
% - Gold-VIX: positive beta during all events → universal safe haven
% - Oil-EM: significant positive beta during Hormuz → oil channel active


\section{Robustness}
\label{sec:dcc_robustness}

% TODO: Rolling-window OLS correlations (60-day, 90-day) as alternative
% to DCC-GARCH → do results hold?
% TODO: Sensitivity to event window definitions
